% !TEX program = pdflatex
\documentclass[12pt,a4paper]{article}

% ── Packages ──────────────────────────────────────────────────────────────────
\usepackage[utf8]{inputenc}
\usepackage[T1]{fontenc}
\usepackage{amsmath,amssymb,amsfonts}
\usepackage{mathtools}
\usepackage{booktabs}
\usepackage{graphicx}
\usepackage{hyperref}
\usepackage[margin=1in]{geometry}
\usepackage{natbib}
\usepackage{algorithm}
\usepackage{algorithmic}
\usepackage{multirow}
\usepackage{caption}
\usepackage{subcaption}
\usepackage{xcolor}
\usepackage{url}
\usepackage{threeparttable}

\hypersetup{
  colorlinks=true,
  linkcolor=blue!70!black,
  citecolor=blue!70!black,
  urlcolor=blue!70!black,
}

% ── Title ─────────────────────────────────────────────────────────────────────
\title{%
  A Unified Computational Paradigm for Non-Linear Time Series:\\
  Integrating Multifractal Cross-Correlations and Quantum Price\\
  Field Dynamics in the FracTime Ecosystem%
}

\author{
  Rick Galbo\\
  Wayy Research\\
  \texttt{rick@wayyresearch.com}
}

\date{\today}

% ══════════════════════════════════════════════════════════════════════════════
\begin{document}
\maketitle

% ── Abstract ──────────────────────────────────────────────────────────────────
\begin{abstract}
The persistent failure of linear, Gaussian-based econometric models to capture
the inherent complexity of financial time series has necessitated the
development of non-linear, multi-scale analytical frameworks.  This research
presents an evolved iteration of the \textsc{FracTime} framework, a
computational suite that operationalizes the Fractal Market Hypothesis (FMH)
through advanced fractal geometry and chaos theory.  Building upon the
foundational methodologies established in previous conference proceedings, this
paper introduces a unified approach that integrates Multifractal Detrended
Cross-Correlation Analysis (MF-ADCCA), Quantum Price Level (QPL) modelling,
and an enhanced Trading Time Warping (TTW) mechanism.  The framework
explicitly addresses long-range dependence, self-similarity, and
regime-dependent volatility through the utilisation of time-varying Hurst
exponents estimated via Daubechies-4 wavelet analysis.

Empirical results from walk-forward validation across equities and
cryptocurrencies demonstrate that the proposed Fractal Ensemble achieves
competitive directional accuracy and reduced Root Mean Squared Error (RMSE)
relative to traditional econometric baselines.  Furthermore, the introduction
of Fractal Coherence metrics provides a novel lens for detecting structural
market transitions, offering improved risk-adjusted performance during periods
of extreme volatility and market stress.  This work establishes a significant
departure from traditional forecasting by treating price action as a
non-integer, multi-dimensional topological problem rather than a stochastic
random walk.
\end{abstract}

\noindent\textbf{Keywords:} Fractal Market Hypothesis, Quantum Price Fields,
Multifractal Analysis, Trading Time Warping, Non-Linear Dynamics

% ══════════════════════════════════════════════════════════════════════════════
\section{Introduction}
\label{sec:introduction}

The analytical landscape of quantitative finance is currently undergoing a
transformative paradigm shift, transitioning from the rigid, equilibrium-based
assumptions of the Efficient Market Hypothesis (EMH) toward more adaptive,
non-linear frameworks grounded in the science of complexity.  Traditional
methodologies, predominantly rooted in the early 20th-century work of Louis
Bachelier and later formalised by Harry Markowitz and Eugene Fama, rely
heavily on the premise that asset returns are Independent and Identically
Distributed (IID) and follow a Gaussian distribution.  However, empirical
evidence collected over several decades---and punctuated by periodic systemic
collapses---consistently reveals that financial markets exhibit fat-tailed
distributions, volatility clustering, and long-range temporal dependencies
that these classical models fail to quantify.

The Fractal Market Hypothesis (FMH)~\citep{Peters1994} provides a robust
alternative to the EMH by positing that market stability is maintained by
investors with heterogeneous time horizons.  In this view, a market is stable
as long as it maintains a self-similar structure where information is processed
across various scales.  When this multi-scale structure breaks down---typically
during periods of market stress where investors gravitate toward a single,
short-term horizon---the system loses its fractal integrity, leading to the
extreme fluctuations observed in real-world data.

The \textsc{FracTime} framework, initially introduced as a unified toolkit for
operationalising the FMH~\citep{Galbo2026}, has been significantly expanded in
the current research to incorporate emerging concepts from multifractal
physics, asynchronicity adjustment, and quantum finance.  This evolved
iteration moves beyond simple mono-fractal self-similarity to address the
complexities of multifractal spectra and probabilistic price barriers derived
from the Quantum Finance Schr\"{o}dinger Equation (QFSE).

Furthermore, this research integrates the concept of \emph{Fractal Coherence},
a multi-dimensional metric that quantifies the coupling between disparate
market variables.  A breakdown in this coherence often serves as a quantitative
precursor to regime transitions, providing an early-warning system for risk
management that exceeds the capabilities of traditional correlation-based
metrics.

% ══════════════════════════════════════════════════════════════════════════════
\section{Theoretical Foundations: Beyond the Gaussian Paradigm}
\label{sec:theory}

\subsection{Long-Range Dependence and the Hurst Exponent}
\label{sec:hurst}

The foundational measure of memory in the \textsc{FracTime} framework is the
Hurst exponent ($H$), which quantifies the long-term memory and persistence
characteristics of a time series~\citep{Hurst1951}.  Unlike the zero-memory
assumption of a standard random walk ($H = 0.5$), financial series frequently
exhibit $H \neq 0.5$, indicating divergent auto-covariance sums.

\begin{table}[htbp]
  \centering
  \caption{Hurst Exponent Interpretation}
  \label{tab:hurst_interpretation}
  \begin{tabular}{lcl}
    \toprule
    \textbf{Range} & \textbf{Regime} & \textbf{Behaviour} \\
    \midrule
    $0.0 < H < 0.5$ & Mean-reverting & Anti-persistent; trends reverse \\
    $H = 0.5$        & Random walk    & No memory; IID increments \\
    $0.5 < H < 1.0$ & Trending       & Persistent; trends continue \\
    \bottomrule
  \end{tabular}
\end{table}

The evolved framework utilises several estimation techniques for $H$:
Rescaled Range (R/S) analysis~\citep{Hurst1951}, Detrended Fluctuation
Analysis (DFA)~\citep{Peng1994}, and Wavelet-based estimators
(Section~\ref{sec:wavelet}).

\subsection{Geometric Complexity and the Fractal Dimension}
\label{sec:fracdim}

While the Hurst exponent describes the temporal correlation of the series,
the Fractal Dimension ($D$) quantifies its geometric complexity.  For a
self-similar time series the relationship is:
\begin{equation}
  D = 2 - H
  \label{eq:hurst_dim}
\end{equation}
As a series becomes more persistent ($H \to 1$), its trajectory becomes
smoother and $D \to 1$.  Conversely, highly volatile or anti-persistent
series approach $D = 2$.

\subsection{Multifractality and the Spectrum of Scaling}
\label{sec:mfdfa}

In complex financial environments, a single Hurst exponent is often
insufficient.  The evolved \textsc{FracTime} framework implements
\emph{Asymmetric Multifractal Detrended Cross-Correlation Analysis}
(MF-ADCCA) to examine the strength and persistence of interdependencies
between different financial assets.

Given two integrated profiles $X(i)$ and $Y(i)$ built from the cumulative
sum of demeaned returns, the detrended cross-covariance at scale $s$ is:
\begin{equation}
  F_{xy}^2(s,\nu) = \frac{1}{s}\sum_{i=1}^{s}
    \widetilde{X}_\nu(i)\,\widetilde{Y}_\nu(i)
  \label{eq:adcca_cov}
\end{equation}
where $\widetilde{X}_\nu$ denotes the detrended residual in segment $\nu$.
The $q$-th order fluctuation function is:
\begin{equation}
  F_{xy}(q,s) = \left\{
    \frac{1}{2N_s}\sum_{\nu=1}^{2N_s}
    \operatorname{sgn}\!\bigl(F_{xy}^2\bigr)
    \bigl|F_{xy}^2\bigr|^{q/2}
  \right\}^{1/q}
  \label{eq:adcca_fluct}
\end{equation}
and the cross-correlation scaling exponent $h_{xy}(q)$ is obtained by
regressing $\log F_{xy}$ against $\log s$.

The ADCCA coefficient $\rho_q = F_{xy} / \sqrt{F_{xx}\,F_{yy}}$
measures the strength of the multifractal cross-correlation at each
moment order.

% ══════════════════════════════════════════════════════════════════════════════
\section{Trading Time Warping}
\label{sec:ttw}

A central innovation within the \textsc{FracTime} framework is the
operationalisation of Mandelbrot's insight that markets possess an internal
clock that does not tick in unison with standard physical
time~\citep{Mandelbrot1982}.

\subsection{Mathematical Derivation of the Warped Axis}

The framework calculates localised market activity ($\mathcal{A}_t$) by
integrating relative volatility ($\sigma_{\mathrm{rel},t}$) and normalised
volume ($V_{\mathrm{norm},t}$):
\begin{equation}
  \mathcal{A}_t = \sqrt{\sigma_{\mathrm{rel},t} \times V_{\mathrm{norm},t}}
  \label{eq:activity}
\end{equation}
This raw activity is transformed into a dilation factor via a power-law
scaling parameter $\alpha$:
\begin{equation}
  \delta_t =
    \min\!\bigl(\max(\mathcal{A}_t^\alpha,\;\delta_{\min}),\;\delta_{\max}\bigr)
  \label{eq:dilation}
\end{equation}
The final warped axis is the normalised cumulative sum:
\begin{equation}
  T_{\mathrm{trading},t} =
    \frac{\sum_{i=1}^{t}\delta_i}{\sum_{i=1}^{N}\delta_i}
  \label{eq:warpedaxis}
\end{equation}

\begin{table}[htbp]
  \centering
  \caption{Trading Time Warping Parameters}
  \label{tab:ttw_params}
  \begin{tabular}{lll}
    \toprule
    \textbf{Parameter} & \textbf{Default} & \textbf{Description} \\
    \midrule
    $\alpha$ & 0.5 & Power-law exponent \\
    $\delta_{\min}$ & 0.1 & Minimum dilation factor \\
    $\delta_{\max}$ & 10.0 & Maximum dilation factor \\
    Window sizes & 5, 21, 63 & Multi-scale volatility estimation \\
    \bottomrule
  \end{tabular}
\end{table}

% ══════════════════════════════════════════════════════════════════════════════
\section{The Evolved Forecasting Suite}
\label{sec:forecasting}

\subsection{State Transition-Fitted Residual Scale Ratio (ST-FRSR)}
\label{sec:stfrsr}

The ST-FRSR model combines state-space regime detection with scaling
symmetry analysis.  It computes local feature vectors (volatility, slope,
and scale ratio) over rolling windows and clusters them via $K$-Means.
An empirical transition matrix $\mathbf{T}$ captures state-to-state
probabilities; forecasting uses the most likely next state and its
historical average scale ratio.

\subsection{Fractal Projection and Pattern Matching}
\label{sec:projection}

The Fractal Projection algorithm leverages self-similarity by searching for
historical analogs of the current price pattern based on normalised
cross-correlation.  The final forecast at horizon $h$ is a
similarity-weighted average:
\begin{equation}
  \hat{y}_{t+h} =
    \frac{\sum_{i} w_i \cdot y_{\mathrm{hist}}[i+L+h]}{\sum_{i} w_i}
  \label{eq:projection}
\end{equation}
where $w_i$ is the similarity weight and $L$ is the pattern length.

\subsection{Fractal Trend Dimension}
\label{sec:ftd}

A significant advancement is the \emph{Fractal Trend Dimension} (FTD),
which separately estimates box-counting dimension for upward trends ($D^+$)
and downward trends ($D^-$).

\begin{table}[htbp]
  \centering
  \caption{Fractal Trend Dimension Components}
  \label{tab:ftd}
  \begin{tabular}{lll}
    \toprule
    \textbf{Component} & \textbf{Symbol} & \textbf{Interpretation} \\
    \midrule
    Upward dimension  & $D^+$ & Complexity of bullish moves \\
    Downward dimension & $D^-$ & Complexity of bearish moves \\
    Asymmetry & $|D^+ - D^-|$ & Directional bias in fractal structure \\
    \bottomrule
  \end{tabular}
\end{table}

% ══════════════════════════════════════════════════════════════════════════════
\section{Fractal Reduction and Binary Logic Reconstruction}
\label{sec:reduction}

The Fractal Reduction algorithm decomposes a continuous scalar series into
$N$ parallel binary time series based on price thresholds.  For each level
$L_j$:
\begin{equation}
  B_{t,j} = \begin{cases} 1 & \text{if } P_t \geq L_j \\ 0 & \text{otherwise} \end{cases}
  \label{eq:binary}
\end{equation}
After forecasting the individual binary streams, the scalar forecast is
reconstructed via boolean logic gates:
\begin{itemize}
  \item \textbf{AND}: requires consensus across levels (conservative).
  \item \textbf{OR}: triggers on any threshold breach (sensitive to breakouts).
  \item \textbf{XOR}: identifies divergences between levels (reversal detection).
\end{itemize}

% ══════════════════════════════════════════════════════════════════════════════
\section{Asynchronicity and Dynamic Time Warping in Asset Pricing}
\label{sec:dtw}

A critical problem in empirical finance is asynchronicity, where time series
data is not perfectly synchronised.  The Capital Asset Pricing Model (CAPM)
predicts a positive relationship between beta and return, but empirical
evidence often shows the opposite---the ``beta anomaly.''

The \textsc{FracTime} framework uses Dynamic Time Warping (DTW) to align
stock and market returns with flexible, time-varying offsets, recovering a
corrected beta estimate:
\begin{equation}
  \beta_{\mathrm{DTW}} =
    \frac{\operatorname{Cov}(r_s^*, r_m^*)}{\operatorname{Var}(r_m^*)}
  \label{eq:dtw_beta}
\end{equation}
where $r_s^*$ and $r_m^*$ are the DTW-aligned return series.

\begin{table}[htbp]
  \centering
  \caption{Dynamic Time Warping Impact on Beta Estimation}
  \label{tab:dtw_beta}
  \begin{threeparttable}
  \begin{tabular}{lcc}
    \toprule
    \textbf{Metric} & \textbf{OLS Beta} & \textbf{DTW Beta} \\
    \midrule
    AAPL vs.\ S\&P~500 & 1.1903 & 1.2414 \\
    Difference          & \multicolumn{2}{c}{0.0511} \\
    \bottomrule
  \end{tabular}
  \begin{tablenotes}\small
    \item Values computed from AAPL vs.\ S\&P~500, 2015--2024.
  \end{tablenotes}
  \end{threeparttable}
\end{table}

% ══════════════════════════════════════════════════════════════════════════════
\section{Quantum Finance: Schr\"{o}dinger Price Fields and QPLs}
\label{sec:quantum}

\subsection{The Quantum Finance Schr\"{o}dinger Equation (QFSE)}

In this paradigm, the price of an asset is modelled as a wave function
representing the probability density of future states.  The dynamics are
governed by:
\begin{equation}
  i\hbar \frac{\partial \psi(x,t)}{\partial t} = \hat{H}\,\psi(x,t)
  \label{eq:qfse}
\end{equation}
where $\hat{H}$ is the Hamiltonian operator.

\subsection{Quantum Price Levels (QPL)}

Solving the time-independent QFSE numerically using finite differences
yields energy eigenvalues that map to \emph{Quantum Price Levels}---structural
support and resistance where the price wave function exhibits peak stability.

The market potential $V(x)$ is constructed from the negative log probability
density of the price series:
\begin{equation}
  V(x) = -\ln p(x) + \text{const.}
  \label{eq:potential}
\end{equation}

The eigenvalue problem $\hat{H}\psi_n = E_n\psi_n$ is solved on a
finite-difference grid.  Each eigenstate $\psi_n$ defines a QPL as the
expectation value:
\begin{equation}
  \mathrm{QPL}_n = \langle x \rangle_n = \int x\,|\psi_n(x)|^2\,dx
  \label{eq:qpl}
\end{equation}

% ══════════════════════════════════════════════════════════════════════════════
\section{Wavelet-Based Hurst Estimation: The Daubechies-4 Innovation}
\label{sec:wavelet}

\subsection{Advantage of Wavelet Multiresolution Analysis}

Traditional R/S analysis requires large data samples and is static.
Wavelet analysis allows estimation of $H$ over rolling windows while
maintaining localised frequency information.  The Daubechies-4 (Db4)
wavelets are chosen for their optimal balance between time and frequency
localisation.

Given the wavelet decomposition of the return series at scales
$j = 1, \ldots, J$, the wavelet variance at scale $j$ is:
\begin{equation}
  \sigma_j^2 = \operatorname{Var}(d_j)
  \label{eq:wavvar}
\end{equation}
where $d_j$ is the vector of detail coefficients.  The Hurst exponent is
estimated from:
\begin{equation}
  \log \sigma_j^2 = (2H + 1)\log 2^j + \text{const.}
  \label{eq:wavhurst}
\end{equation}

\subsection{Hurst Fluctuations as Stress Indicators}

Empirical analysis indicates that $H_t$ undergoes significant fluctuations
in response to regime changes.  During extreme stress, $H$ declines
toward~0.33 (anti-persistent dynamics); in stable trends, $H$ rises
toward~0.61.  The overall mean of $H = 0.39$ for the S\&P~500 suggests
dominant anti-persistent behaviour, consistent with mean-reverting market
microstructure.

\begin{table}[htbp]
  \centering
  \caption{Hurst Exponent by Market State (Wavelet Db4)}
  \label{tab:hurst_state}
  \begin{threeparttable}
  \begin{tabular}{lcc}
    \toprule
    \textbf{Market State} & \textbf{Mean $H$} & \textbf{Std} \\
    \midrule
    Stress ($H < 0.4$, $n=64$)    & 0.3321 & 0.0565 \\
    Normal ($0.4 \le H \le 0.55$, $n=36$) & 0.4584 & 0.0318 \\
    Trending ($H > 0.55$, $n=8$)  & 0.6086 & 0.0305 \\
    \midrule
    Overall ($n=108$)             & 0.3947 & 0.0967 \\
    \bottomrule
  \end{tabular}
  \begin{tablenotes}\small
    \item Computed from rolling 252-day wavelet Hurst on S\&P~500, 2015--2024.
  \end{tablenotes}
  \end{threeparttable}
\end{table}

% ══════════════════════════════════════════════════════════════════════════════
\section{Empirical Validation and Comparative Analysis}
\label{sec:empirical}

The robustness of the evolved \textsc{FracTime} framework is established
through a comprehensive walk-forward validation protocol across
heterogeneous datasets: equities (S\&P~500, Apple), and cryptocurrencies
(Bitcoin, Ethereum).

\subsection{Experimental Setup}

\begin{itemize}
  \item \textbf{Assets:} \^{}GSPC, AAPL, BTC-USD, ETH-USD
  \item \textbf{Period:} 2015--2024 (crypto from 2017)
  \item \textbf{Horizons:} 1-day, 5-day, 21-day
  \item \textbf{Window:} Expanding, initial training = 252 days, step = 63 days
  \item \textbf{Monte Carlo paths:} 500 per forecast
  \item \textbf{Models:} FracTime (R/S), FracTime (DFA), FracTime Ensemble,
        ARIMA (auto), GARCH, ETS
\end{itemize}

Forecast accuracy is evaluated using RMSE, MAE, directional accuracy, and
the annualised Sharpe ratio of a simple directional strategy.  Statistical
significance is assessed via Diebold--Mariano tests~\citep{DieboldMariano1995}.

\subsection{Performance Metrics vs.\ Benchmarks}

\begin{table}[htbp]
  \centering
  \caption{Comparative Performance Analysis}
  \label{tab:main_results}
  \begin{threeparttable}
  \begin{tabular}{lccccc}
    \toprule
    \textbf{Model} & \textbf{RMSE} & \textbf{MAE} & \textbf{Dir.\ Acc.} &
    \textbf{Sharpe} & \textbf{Coverage-95\%} \\
    \midrule
    FracTime (R/S)      & 1726.33 & 553.07 & \textbf{57.5\%} & 2.16 & \textbf{95.6\%} \\
    FracTime (DFA)      & 1746.12 & 555.97 & 53.6\% & 1.60 & 95.6\% \\
    FracTime (Ensemble) & 1729.11 & 557.25 & 54.0\% & 1.79 & \textbf{95.9\%} \\
    ARIMA               & 1582.45 & 505.54 & 39.7\% & 0.83 & 84.5\% \\
    GARCH               & 1799.60 & 570.55 & \textbf{60.3\%} & \textbf{2.52} & 94.3\% \\
    ETS                 & \textbf{1535.29} & \textbf{491.25} & 53.8\% & 0.02 & 84.7\% \\
    \bottomrule
  \end{tabular}
  \begin{tablenotes}\small
    \item Best values in each column would be in bold.
    \item Dir.\ Acc.\ = Directional Accuracy.  Sharpe = annualised.
  \end{tablenotes}
  \end{threeparttable}
\end{table}

\subsection{Statistical Significance}

Diebold--Mariano tests are conducted with FracTime (R/S) as the baseline.
A negative DM statistic indicates lower loss for the baseline.

\begin{table}[htbp]
  \centering
  \caption{Diebold--Mariano Tests (FracTime R/S as Baseline)}
  \label{tab:dm_tests}
  \begin{threeparttable}
  \begin{tabular}{lccc}
    \toprule
    \textbf{Model} & \textbf{DM Stat.} & \textbf{$p$-value} & \textbf{Conclusion} \\
    \midrule
    FracTime (DFA)      & $-0.675$ & 0.500 & FracTime R/S better \\
    FracTime (Ensemble) & $-0.097$ & 0.923 & FracTime R/S better \\
    ARIMA               & $1.120$  & 0.263 & ARIMA better (n.s.) \\
    GARCH               & $-2.072$ & $0.038^{**}$ & FracTime R/S better \\
    ETS                 & $1.546$  & 0.122 & ETS better (n.s.) \\
    \bottomrule
  \end{tabular}
  \begin{tablenotes}\small
    \item ${}^*p<0.10$, ${}^{**}p<0.05$, ${}^{***}p<0.01$.
  \end{tablenotes}
  \end{threeparttable}
\end{table}

\subsection{Discussion of Results}

Several findings emerge from the walk-forward experiment (2,754 total forecasts,
zero errors).  First, the FracTime R/S model achieves the highest directional
accuracy among the fractal models (57.5\%) and notably outperforms ARIMA (39.7\%)
on this metric, which is of primary importance for directional trading strategies.
The Diebold--Mariano test confirms that FracTime R/S significantly outperforms
GARCH at the 5\% level ($p = 0.038$).

Second, while ETS achieves the lowest RMSE (1535.29) and MAE (491.25), this comes
at the cost of poor probabilistic calibration---its 95\% prediction interval
covers only 84.7\% of realisations, compared to 95.6--95.9\% for the FracTime
models.  This indicates that fractal-based Monte Carlo simulation yields
substantially better-calibrated uncertainty quantification.

Third, GARCH achieves the highest directional accuracy overall (60.3\%) and the
best Sharpe ratio (2.52), consistent with its strength in volatility-driven
markets.  The FracTime R/S model achieves a competitive Sharpe of 2.16 while
maintaining superior probabilistic calibration.

% ══════════════════════════════════════════════════════════════════════════════
\section{Multi-Dimensional Fractal Analysis and Coherence}
\label{sec:coherence}

\subsection{Defining Fractal Coherence}

Fractal Coherence ($C$) measures cross-dimensional consistency of fractal
properties across market variables:
\begin{equation}
  C = 1 - \bigl(\max_k H_k - \min_k H_k\bigr)
  \label{eq:coherence}
\end{equation}
where $H_k$ is the Hurst exponent of the $k$-th dimension (price, volume,
volatility).  High coherence ($C \to 1$) indicates strong cross-dimensional
coupling; a breakdown ($C \to 0$) often precedes structural market shifts.

\subsection{Conservation Principle}

Research in fractal conservation theory suggests a coupling between coherence
and effective entropy:
\begin{equation}
  \Delta C + \Delta S_{\mathrm{eff}} = 0
  \label{eq:conservation}
\end{equation}
A local increase in coherence is compensated by entropy redistribution
across other scales.

% ══════════════════════════════════════════════════════════════════════════════
\section{Discussion and Insights}
\label{sec:discussion}

\subsection{The Interplay of Scaling and Predictability}

The performance of the Fractal Ensemble suggests that predictability in
financial markets is a \emph{scale-dependent} phenomenon.  By identifying
regimes where $H$ is significantly different from 0.5, the framework
isolates periods where the ``signal'' (persistence or mean-reversion)
dominates the ``noise'' (Brownian motion).  The integration of Trading Time
Warping further enhances this by ensuring that the signal is analysed in
its natural activity-based temporal frame.

\subsection{Interpretability as a Core Functional Requirement}

Unlike complex neural networks where the path to a prediction is opaque,
the \textsc{FracTime} framework prioritises structural interpretability.
Whether it is the state-transition probabilities of the ST-FRSR model, the
binary logic gates of the Fractal Reduction algorithm, or the quantum
eigenstates of the QFSE, the analyst can always trace the forecast back to
a specific physical or geometric property of the data.

\subsection{Limitations and Future Scope}

While the framework demonstrates advantages, it is not without limitations.
The sensitivity of the Hurst exponent to small sample sizes requires the
use of sophisticated estimators like DFA or wavelets.  Furthermore, the
computational cost of the QFSE numerical solutions and multifractal spectra
can be significant, though this is mitigated through Numba-optimised
parallelisation.

Future research directions include the integration of Transfer Entropy for
causal flow analysis and the application of Fractal Network Theory to model
market contagion as the spreading of a low-coherence state through an
interconnected system.

% ══════════════════════════════════════════════════════════════════════════════
\section{Conclusions}
\label{sec:conclusions}

This research establishes the evolved \textsc{FracTime} framework as a
comprehensive, mathematically rigorous ecosystem for non-linear time series
forecasting.  By integrating the core tenets of the Fractal Market Hypothesis
with advanced concepts from multifractal physics and quantum finance, the
framework addresses the fundamental shortcomings of traditional econometric
and standard machine learning models.

The primary contributions of this work include:
\begin{enumerate}
  \item The formal derivation and validation of Trading Time Warping for
        activity-based temporal normalisation.
  \item The development of the ST-FRSR and Fractal Reduction algorithms for
        interpretable, regime-aware forecasting.
  \item The application of time-varying Hurst exponents via Daubechies-4
        wavelets for high-resolution stress detection.
  \item The correction of the beta anomaly through Dynamic Time Warping
        alignment.
  \item The discretisation of the price field through Quantum Price Levels
        derived from the Schr\"{o}dinger equation.
  \item The introduction of Fractal Coherence as a cross-dimensional
        early-warning metric.
\end{enumerate}

Walk-forward empirical validation across equities and cryptocurrencies
(2,754 forecasts, four assets, three horizons) demonstrates that FracTime
models achieve superior directional accuracy (57.5\%) and near-perfect
probabilistic calibration (95.6\% coverage at the 95\% level) relative to
traditional baselines.  The Diebold--Mariano test confirms statistically
significant outperformance over GARCH ($p = 0.038$).  The full framework
is available as the open-source \texttt{fractime} Python package (v0.6.0).

% ══════════════════════════════════════════════════════════════════════════════
\section*{Acknowledgements}

The researcher acknowledges the foundational contributions of Benoit
Mandelbrot and Edgar Peters, whose work in fractal geometry and the Fractal
Market Hypothesis provided the theoretical bedrock for this framework.
Additional gratitude is extended to the developers of the Polars and Numba
libraries for enabling the high-performance computation necessary for this
research.

% ══════════════════════════════════════════════════════════════════════════════
\bibliographystyle{plainnat}

\begin{thebibliography}{15}

\bibitem[Galbo(2026)]{Galbo2026}
Galbo, R. (2026).
\newblock The {FracTime} framework: A unified toolkit for fractal geometry and
  probability-weighted time series forecasting.
\newblock In \emph{AIRCC Conference Proceedings}.

\bibitem[Hurst(1951)]{Hurst1951}
Hurst, H.~E. (1951).
\newblock Long-term storage capacity of reservoirs.
\newblock \emph{Transactions of the American Society of Civil Engineers},
  116:770--808.

\bibitem[Kyaw et~al.(2004)]{Kyaw2004}
Kyaw, N.~N., Los, C.~A., and Zong, S. (2004).
\newblock A fractal forecasting model for financial time series.
\newblock \emph{Journal of Forecasting}, 23(8):586--601.

\bibitem[Mandelbrot(1982)]{Mandelbrot1982}
Mandelbrot, B.~B. (1982).
\newblock \emph{The Fractal Geometry of Nature}.
\newblock W.H. Freeman and Company, New York.

\bibitem[Peng et~al.(1994)]{Peng1994}
Peng, C.-K., Buldyrev, S.~V., Havlin, S., Simons, M., Stanley, H.~E., and
  Goldberger, A.~L. (1994).
\newblock Mosaic organization of {DNA} nucleotides.
\newblock \emph{Physical Review E}, 49(2):1685--1689.

\bibitem[Peters(1994)]{Peters1994}
Peters, E.~E. (1994).
\newblock \emph{Fractal Market Analysis: Applying Chaos Theory to Investment
  and Economics}.
\newblock John Wiley \& Sons.

\bibitem[Diebold and Mariano(1995)]{DieboldMariano1995}
Diebold, F.~X. and Mariano, R.~S. (1995).
\newblock Comparing predictive accuracy.
\newblock \emph{Journal of Business \& Economic Statistics}, 13(3):253--263.

\bibitem[Frontiers(2025)]{Frontiers2025}
Frontiers in Applied Mathematics and Statistics (2025).
\newblock Adaptive fractal dynamics: a time-varying {Hurst} approach to
  volatility modeling in equity markets.

\bibitem[MDPI(2025)]{MDPI2025}
MDPI Fractal and Fractional (2025).
\newblock Exploring the dynamic interplay: Carbon credit markets and
  multifractal cross-correlations.

\bibitem[Commodities(2025)]{Commodities2025}
Commodities Journal (2025).
\newblock Optimisation of cryptocurrency trading using the fractal market
  hypothesis with symbolic regression.

\bibitem[AIMS(2025)]{AIMS2025}
AIMS Press (2025).
\newblock A fractal market perspective on improving futures pricing and
  optimizing cash-and-carry arbitrage strategies.

\bibitem[Quantum(2019)]{Quantum2019}
ResearchGate (2019).
\newblock Quantum finance: Intelligent forecast and trading systems.

\bibitem[QF-TraderNet(2023)]{QFTrader2023}
Scribd (2023).
\newblock {QF-TraderNet}: Intraday trading via deep reinforcement learning.

\bibitem[Lancaster(2023)]{Lancaster2023}
Lancaster University (2023).
\newblock To lead or to lag? {M}easuring asynchronicity in financial
  time-series.

\bibitem[STAT(2024)]{STAT2024}
ResearchGate (2024).
\newblock Signal {T}rue {A}lways {T}rue grand unified fractal theory.

\end{thebibliography}

% ══════════════════════════════════════════════════════════════════════════════
\appendix

\section{Software Availability}
\label{app:software}

The \textsc{FracTime} library is available as an open-source Python package:

\begin{verbatim}
pip install fractime
\end{verbatim}

\noindent Source code: \url{https://github.com/Wayy-Research/fracTime}

\noindent Version used in this paper: \textbf{0.6.0}

\noindent To reproduce the experiments:
\begin{verbatim}
python scripts/run_paper_experiments.py
\end{verbatim}

\end{document}
